\chapter{Обзор методов дифференцируемого рендеринга поверхностей}
\label{ch:chap1}

Понятие дифференцируемого рендеринга, решающего задачу реконструкции
параметров сцены по её изображениям
с помощью градиентных методов оптимизации, было введено в 2014 году
в проекте OpenDR \cite{OpenDR2014}. В этой же работе была представлена
первая открытая одноимённая
система для приближённого решения поставленной задачи: дифференцируемый
рендерер, позволявший сделать отрисовку сцены и получить градиенты относительно
её параметров. С тех пор данное направление активно развивается, предлагая
большое разнообразие решений.

На данный момент можно выделить несколько групп методов дифференцируемого
рендеринга, различающихся типом используемого представления для реконструкции
объектов сцены, а также методом их отрисовки. Первая группа:
методы, использующие объёмные представления.
К ней относятся такие методы, как NeRF \cite{NeRF2020} и 3D gaussian splatting (3DGS) \cite{3DGS}.
Преимущество этих методов заключается в том, что алгоритм их рендеринга полностью
дифференцируемый, благодаря чему они восстанавливают сцену без специфических
оценок, корректирующих градиенты.
Основным недостатком является то, что эти представления не могут быть использованы
в физически-обоснованном рендеринге напрямую: для этого их нужно конвертировать в
поверхностные представления, например полигональный меш. Подобная конвертация
является нетривиальной задачей, также при конвертации неизбежны потери в
точности результата реконструкции.

Вторая группа использует объёмный рендеринг для восстановления
поверхностных представлений. Сюда относятся:
\begin{enumerate}
    \item Методы, восстанавливающие поверхностные представления на основе идеи 3DGS:
          2D gaussian splatting \cite{2DGS} заменяет объёмные гауссианы на
          плоские диски, нормали которых совпадают с нормалями поверхностей объектов
          в сцене. SuGaR \cite{SuGaR} применяет регуляризацию, чтобы 3D-
          гауссианы выравнивались вдоль восстанавливаемой поверхности, после чего
          получает меш. В прошлом году также появились методы, направленные
          на реконструкцию меша с помощью алгоритма сплаттинга: MiLO \cite{MiLO}
          дифференцируемым образом получает меш из результата реконструкции с
          помощью 3DGS. Triangle splatting \cite{TrigSplat} заменяет гауссианы на
          треугольники в алгоритме реконструкции, а Mesh splatting \cite{MeshSplat}
          обеспечивает связность примитивов, получая в результате восстановления
          полноценный меш.


    \item Нейронные методы, восстанавливающие функцию расстояния со знаком (SDF) --
          неявное представление поверхности, представляющее собой скалярное поле
          дистанций до поверхности в пространстве сцены. NeuS \cite{NeuS}
          предложил новую функцию взвешивания для восстановления SDF с помощью NeRF,
          что позволяет оптимизировать SDF-представление
          без маски объекта и обеспечивает качественную реконструкцию объектов
          сложной структуры. NeuS2 \cite{NeuS2} заменяет MLP-кодировщик позиции на
          хеш-таблицы \cite{InstantNGP}, достигая существенного ускорения обучения
          (несколько минут вместо нескольких часов) без потери точности восстановленных
          объектов. Метод также можно применять для динамических сцен. NeuralAngelo \cite{NeuralAngelo}
          в целях повышения точности использует численные градиенты более высокого
          порядка для сглаживания нормалей.
\end{enumerate}
Проблемы данных методов -- необходимость выбирать между скоростью и качеством
реконструкции, а также проблема, свойственная и первой группе: либо восстановленные
модели требуется дополнительно конвертировать в традиционные представления,
либо они вовсе хранятся в весах нейронной сети; процесс конвертации является
нетривиальной задачей.

В третью группу методов можно отнести системы, реализующие дифференцируемую
версию классического алгоритма растеризации. NVDiffRast \cite{NVDiffRast}
предлагает модульный дифференцируемый рендерер на базе растеризации, чьи
базовые операции -- растеризация, интерполяция атрибутов, фильтрация текстур
и антиалиасинг -- реализованы как высокопроизводительные CUDA‑ядра, которые
можно интегрировать непосредственно в фреймворки автоматического дифференцирования.
NVDiffRast намеренно не фиксирует модель камеры или освещения, предоставляя лишь
низкоуровневые строительные блоки графического конвейера, поверх которых
можно свободно реализовывать собственные схемы затенения и оптимизации.
Подобная модульность и ориентация на отложенное затенение позволяют достигать
высокой производительности и численной стабильности градиентов при рендеринге
сцен с большим количеством треугольников в высоком разрешении.
На базе NVDiffRast построен метод NVDiffRec \cite{NVDiffRec}, оптимизирущий
треугольные меши с материалами и освещением, что обеспечивает совместимость
результата реконструкции с традиционными растеризаторами и трассировщиками
лучей. Геометрия параметризуется с помощью SDF и нейронного представления
DMTet, тогда как материалы задаются через 3D сетки, а освещение моделируется
в виде HDR‑окружения. Тем самым nvdiffrec демонстрирует, что комбинация
дифференцируемого рендерера на базе растеризации и параметризаций позволяет
устойчиво решать задачу реконструкции с точностью, достаточной для последующей
физически корректной визуализации.
NVDiffRecMC \cite{NVDiffRecMC} интегрирует в пайплайн nvdiffrec более реалистичную
модель освещения, основанную на дифференцируемой трассировке лучей и выборке
по значимости. Авторы показывают, что использование простых моделей освещения
(предварительно отфильтрованное прямое освещение) ограничивает качество
декомпозиции сцены на форму, материалы и свет, особенно в условиях сложных
взаимных отражений и сложных конфигураций освещения. В NVDiffRecMC сохраняется
параметризация сцены из NVDiffRec, однако процесс прямого рендеринга заменяется
на оценку интеграла освещённости методом Монте-Карло, дополненную дифференцируемым
денойзером, который встраивается в вычислительный граф и существенно снижает
дисперсию как изображений, так и соответствующих градиентов. За счёт этого
становится возможным выполнять градиентную оптимизацию при относительно малом
числе выборок на пиксель.

Альтернативный подход к дифференцируемой растеризации предложен в 2024 году \cite{SGORast}.
Вместо разработки специализированного дифференцируемого рендерера, авторы
рассматривают произвольный недифференцируемый растеризатор как "чёрный ящик"
и оборачивают его в схему стохастической оценки градиентов на основе случайных
возмущений параметров сцены. Классический подход плохо масштабируется по числу
параметров, поскольку требует оценивать градиенты в высокомерном пространстве;
ключевое наблюдение авторов заключается в том, что число параметров, реально
влияющих на один пиксель, ограничено, поэтому можно оценивать градиент попиксельно,
усредняя стохастические оценки по независимым возмущениям только тех параметров,
которые влияют на данный пиксель. Такая локализация делает метод масштабируемым,
позволяя с минимальными усилиями превратить существующий растеризатор (например,
в игровом движке) в дифференцируемый, без зависимости от сторонних библиотек и
без существенных алгоритмических изменений. В отличие от NVDiffRast, NVDiffRec
и NVDiffRecMC, где градиенты либо выводятся аналитически, либо аккуратно пропускаются
через фреймворк автодифференцирования, данный метод опирается на стохастические
оценки, по природе более шумные, но взамен дающие максимальную гибкость и
совместимость с существующей инфраструктурой рендеринга.

Основные ограничения этих методов связаны с обработкой видимости и дискретностью
в растеризации, а также с компромиссом между физической корректностью и временем
работы. NVDiffRast не способен обрабатывать сложные случаи перекрытия геометрии,
также он требует антиалисаинга для улучшения оценки градиентов. NVDiffRec использует
простую модель освещения, а конвертация внутреннего представления в меш сильно
портит качество модели. NVDiffRecMC зашумлён и активно полагается на дифференцируемый
денойзинг, который усложняет и замедляет работу метода. Зашумлённость градиентов
является основной проблемой \cite{SGORast}. Как и для классического метода конечных
разностей, оценки градиента остаются существенно более шумными и менее эффективными,
чем корректные градиенты, полученные через обратное распространение ошибки в
дифференцируемом рендерере. Авторы отмечают, что для снижения дисперсии
требуется усреднение по нескольким стохастическим оценкам, что увеличивает
число прогонов растеризатора на один шаг оптимизации.

Четвёртая группа использует поверхностные представления совместно с трассировкой путей. Основной
проблемой этой группы является необходимость учитывать разрывы, представляющие
собой границы объекта на изображении. Первые решения, корректно учитывающие
их в рамках физически-обоснованного рендеринга без использования нейронных
сетей: Edge Sampling \cite{ES2018}, репараметризация \cite{Reparam2019}, Warped Area Sampling \cite{WAS2020}.
Для корректного расчёта градиентов эти подходы разными способами оценивают
интеграл по разрывам на изображении, называемый граничным интегралом.

Edge sampling \cite{ES2018} напрямую строит оценку граничного интеграла для меша по
рёбрам, принадлежащим силуэту для каждого ракурса, а также внутренним разрывам
видимости. Для этого вводятся специальные структуры данных и процедуры выборки
рёбер, обеспечивающие эффективное нахождение тех участков геометрии, где
небольшое изменение параметра вызывает резкий скачок видимости. В результате
метод даёт оценки градиента практически без смещения для геометрии, освещения
и параметров камеры, корректно учитывает тени и глобальное освещение, но
страдает от высокой дисперсии и ухудшения масштабируемости при сложной
геометрии и множестве силуэтов, особенно в отражениях и преломлениях.

Репараметризация \cite{Reparam2019} вместо явного семплирования границы строит для каждого
рассматриваемого интеграла локальное преобразование переменных, уменьшающее
зависимость положения разрывов от дифференцируемых параметров. Идея состоит
в локальной перепараметризации путей: разрыв видимости "распрямляется" в
координатах интегрирования, так что результат становится более гладкой
функцией параметров, позволяющей применять стандартную стохастическую оценку
и автоматическое дифференцирование. Метод не требует явной выборки рёбер меша
и показывает низкое смещение и разброс в задачах с глянцевыми отражениями и
мягкими тенями, однако построение параметризации происходит "на лету" и
глобально влияет на схему выборки, что усложняет реализацию и может приводить
к росту разброса при сложной конфигурации сцены.

Метод Warped area sampling \cite{WAS2020} предлагает перейти от граничного
интеграла к интегралу по области, используя теорему Остроградского-Гаусса,
и тем самым полностью отказаться от явного семплирования граничного интеграла.
Производная по параметру записывается как интеграл от некоторого специальным
образом сконструированного векторного поля, согласованного с деформацией
границы, после чего формируется оценка методом Монте‑Карло для этого интеграла.
Для практического применения важно, что метод легко интегрируется в
классический прямой трассировщик путей, не требует силуэтных структур и
допускает строго несмещённые оценки при невысокой дисперсии. По сравнению с
\cite{ES2018}, подход \cite{WAS2020} уменьшает затраты времени и повышает
устойчивость реконструкции к сложности геометрии, но платит за это ростом
разброса градиента.

Три вышеописанных метода предлагают алгоритмы, на практике использующие меш
для восстановления объектов, однако это представление плохо работает в задаче
дифференцируемого рендеринга \cite{Nicolet}. Современные методы
реконструкции поверхности, принадлежащие этой группе, проводят реконструкцию
для функций расстояния со знаком. Они включают в себя как адаптацию предыдущих
методов, в частности warped area sampling \cite{WAS2020}, под SDF -- работа
\cite{WASSDF2022}, а также версия для нейронного рендеринга \cite{DSDF2022} -- так и
новые подходы, использующие свойства функций расстояния.

Работа \cite{Relax2024} предлагает метод релаксированной границы для оценки
граничного интеграла. Этот алгоритм для дифференцируемого рендеринга
поверхностей, заданных SDF, гораздо проще в сравнении с \cite{WASSDF2022}.
Авторы подчёркивают, что многие существующие методы дифференцируемого
физически-обоснованного рендеринга требуют либо сложных структур данных для
видимости, либо репараметризаций с глобальным влиянием на дисперсию; цель \cite{Relax2024} --
минимизировать алгоритмическую сложность, сохранив корректную обработку
видимости с использованием SDF. Предлагаемый алгоритм использует свойства
SDF и аккуратную локальную обработку видимости, что даёт сопоставимые или
лучшие градиенты при более прямолинейной реализации и облегчённой интеграции
в задачи с градиентной оптимизацией.

В качестве базового был выбран подход, предлагающий метод релаксированной
границы \cite{Relax2024}, по следующим причинам:
он работает в рамках физически-обоснованного рендеринга; он использует
функции расстояния со знаком, которые в задаче реконструкции поверхности
проявили себя лучше в сравнении с мешами; он предлагает алгоритмически
простой способ корректной оценки градиентов, что обеспечивает быстрое и
точное восстановление поверхности объектов.


\endinput